\begin{table}[H]
\centering
\caption{Dispersión y masa alrededor del salario mínimo por logro educativo, 2010 y 2025}
\label{tab:masa_smlmv_educacion}
\begingroup
\setlength{\tabcolsep}{3.5pt}
\scriptsize
\begin{tabular}{@{}p{3.5cm}rrrrrr@{}}
\toprule
& \multicolumn{2}{c}{P90/P10 mensual} & \multicolumn{2}{c}{P90/P10 por hora} & \multicolumn{2}{c}{Cerca del SMLMV} \\
\cmidrule(lr){2-3} \cmidrule(lr){4-5} \cmidrule(l){6-7}
Logro educativo & 2010 & 2025 & 2010 & 2025 & 2010 & 2025 \\
\midrule
Primaria o menos & 9,2 & 7,2 & 6,3 & 5,2 & 14,2\% & 16,2\% \\
Secundaria o media & 8,3 & 5,0 & 5,4 & 4,0 & 21,0\% & 29,1\% \\
Superior o universitaria & 10,0 & 6,8 & 8,3 & 6,6 & 11,7\% & 18,4\% \\
\midrule
\textbf{Total} & 12,2 & 8,8 & 9,0 & 7,3 & 16,4\% & 22,8\% \\
\bottomrule
\end{tabular}
\endgroup
\caption*{\footnotesize Nota: cerca del SMLMV corresponde a la proporción de ocupados cuya remuneración mensual equivalente está entre 0,9 y 1,1 veces el salario mínimo legal mensual vigente del año, ambos expresados en pesos de 2025. Cálculos ponderados con factores de expansión. Fuente: cálculos propios con GEIH del DANE y decretos del Gobierno nacional para SMLMV.}
\end{table}
